% ============================================================================
% Fichier  : partie4/P04_C18_continuite.tex
% Titre    : Stratégies et Plans de Continuité
% Auteur   : MINKA MI NGUIDJOÏ Thierry Emmanuel
% Version  : 1.0 -- novembre 2026
% ============================================================================
% Sources : ISO 22301:2019 (8.3, 8.4) ; ISO 22313:2020. Opérations GAROUA et
% MOUNGO : cas fictifs, chiffres cohérents entre eux.
% ============================================================================

\chapter{Stratégies et Plans de Continuité}
\label{chap:pca}
\chaptermark{Stratégies et plans de continuité}

\epigraph{%
    « Le plan de continuité se juge le premier matin,\\
    quand ceux qui doivent l'appliquer le trouvent sans le chercher. »%
}{Principe du cours}

\niveauChapitre{Intermédiaire}{Chapitre~\ref{chap:bia}.}

\begin{boxobjectifs}
\begin{itemize}
    \item Situer le système de management de la continuité de
          l'ISO~22301.
    \item Choisir une stratégie de continuité pour chaque type de
          ressource.
    \item Concevoir un mode dégradé.
    \item Rédiger un plan de continuité d'activité~: déclenchement, rôles,
          procédures, ressources.
    \item Articuler plan de gestion de crise, plan de continuité et plan de
          reprise informatique.
\end{itemize}
\end{boxobjectifs}

\bigskip

% ============================================================================
\section*{Question directrice : comment LiZbiZ a-t-elle livré pendant l'inondation de Bonabéri~?}
% ============================================================================

Dans la nuit du 14 au 15~août 2026, l'entrepôt de Bonabéri, à Douala, est
inondé par 40~cm d'eau. Cinq mois plus tôt, une telle nuit aurait arrêté
les expéditions de Douala pour une semaine. Cette fois, LiZbiZ dispose
d'un plan de continuité, écrit en juin après la crise de mars. Les
expéditions reprennent en 36~heures depuis un entrepôt de repli. Ce
chapitre montre comment un tel plan se construit.

\medskip
\textbf{Plan du chapitre.}
\begin{enumerate}
    \item Le système de management de la continuité (section~\ref{sec:c18-smca}).
    \item Les stratégies par ressource (section~\ref{sec:c18-strategies}).
    \item Le mode dégradé (section~\ref{sec:c18-degrade}).
    \item Le plan de continuité (section~\ref{sec:c18-plan}).
    \item Les trois plans (section~\ref{sec:c18-plans}).
\end{enumerate}

\separateur

% ============================================================================
\section{Le Système de Management de la Continuité}
\label{sec:c18-smca}
% ============================================================================

L'\norme{ISO~22301:2019} fixe les exigences d'un système de management de
la continuité d'activité, certifiable~\cite{iso22301_2019}~; l'ISO~22313
en donne les lignes directrices~\cite{iso22313_2020}. Le cœur opérationnel
est la clause~8~: analyse d'impact et appréciation des risques (8.2),
stratégies et solutions (8.3), plans et procédures (8.4), programme
d'exercices (8.5), évaluation de la documentation et des capacités (8.6).

\begin{boxdef}
\termecle{Continuité d'activité}~: capacité d'une organisation à
continuer à fournir des produits et des services dans des délais et à une
capacité acceptables, prédéfinis, pendant une perturbation.
\end{boxdef}

% ============================================================================
\section{Les Stratégies par Ressource}
\label{sec:c18-strategies}
% ============================================================================

\begin{tableaucours}{L{2.6cm}L{4.8cm}L{4.8cm}}{%
    \tableauHeader{Ressource} &
    \tableauHeader{Stratégies possibles} &
    \tableauHeader{Choix de LiZbiZ}
}
Personnes & Polyvalence, suppléants désignés, travail à distance & Deux
suppléants formés par poste clé \\
Locaux & Site de repli, accord de réciprocité, location à la demande &
Accord avec un entrepôt partenaire à Bassa (Douala) \\
Technologies & Site de secours, nuage, matériel en réserve & Reprise du
système central à Yaoundé (Chapitre~\ref{chap:pra}) \\
Information & Sauvegardes hors ligne, documents papier & Sauvegardes
immuables, bordereaux papier numérotés \\
Fournisseurs & Second fournisseur, clauses de continuité, stocks & Second
opérateur à Garoua, transporteurs de réserve \\
\end{tableaucours}
\vspace{-0.8em}
\begin{center}{\small\itshape Tableau~18.1. Stratégies de continuité par ressource}\end{center}

Le choix se fait en comparant, pour chaque stratégie, son \emph{coût} et le
RTO qu'elle permet d'atteindre, au regard du RTO exigé par l'analyse
d'impact. On ne finance pas une stratégie plus rapide que nécessaire.

% ============================================================================
\section{Le Mode Dégradé}
\label{sec:c18-degrade}
% ============================================================================

\begin{boxdef}
\termecle{Mode dégradé}~: manière de poursuivre une activité, à un niveau
réduit mais acceptable, sans une ou plusieurs de ses ressources habituelles.
\end{boxdef}

\begin{tableaucours}{L{3.2cm}L{9.0cm}}{%
    \tableauHeader{Activité} &
    \tableauHeader{Mode dégradé de LiZbiZ}
}
Expéditions & Bordereaux papier numérotés en carnets~; liste des clients
prioritaires~; tournées fixes \\
Suivi & Tableur partagé hors du système principal, mis à jour toutes les
deux heures~; ligne téléphonique dédiée \\
Facturation & Factures reconstituées à partir des bordereaux, émises dès le
retour du système, par ordre d'ancienneté \\
Paiements & Virements urgents seulement, double signature sur papier,
plafond \\
\end{tableaucours}
\vspace{-0.8em}
\begin{center}{\small\itshape Tableau~18.2. Modes dégradés de LiZbiZ}\end{center}

\begin{boxalert}
\textbf{Violation typique.} Un mode dégradé qui dépend de ce qui est en
panne (un tableur stocké sur le serveur chiffré).

\smallskip
\textbf{Règle~:} le mode dégradé utilise des ressources \emph{indépendantes}
de celles qu'il remplace~: papier, autre site, autre fournisseur. Il est
préparé à l'avance (carnets imprimés, listes à jour) et entraîné.
\end{boxalert}

% ============================================================================
\section{Le Plan de Continuité}
\label{sec:c18-plan}
% ============================================================================

\begin{boxprocedure}
\textbf{Structure du plan de continuité d'une activité}
\begin{enumerate}
    \item \textbf{Objet et périmètre}~: activité, sites, objectifs de
          reprise.
    \item \textbf{Critères de déclenchement}~: qui décide, à partir de quel
          seuil.
    \item \textbf{Rôles}~: responsable du plan, suppléants, équipes.
    \item \textbf{Procédures}~: premières heures, mode dégradé, retour à la
          normale.
    \item \textbf{Ressources}~: site de repli, matériel, documents, accès.
    \item \textbf{Contacts}~: internes, clients prioritaires, fournisseurs,
          autorités, assureur.
    \item \textbf{Fin de crise}~: critères de retour à la normale.
\end{enumerate}
\end{boxprocedure}

\begin{boxscenario}{Opération MOUNGO. La nuit du 14 août}
\textbf{Constatations.}
\begin{itemize}
    \item 15~août, 4~h~10~: le gardien signale 40~cm d'eau dans
          l'entrepôt de Bonabéri.
    \item 5~h~30~: M.~NGUELE, responsable du plan, le déclenche par
          téléphone.
    \item 6~h~00~: les équipes appliquent la procédure~: inventaire des
          1\,200~colis, transfert des colis secs, information des clients
          prioritaires.
    \item 15~août, 18~h~00~: l'entrepôt partenaire de Bassa reçoit les
          premiers colis~; 16~août, 16~h~00~: les expéditions de Douala
          reprennent depuis Bassa.
    \item Bilan~: 310~colis endommagés, d'une valeur déclarée de
          96~millions de FCFA~; expéditions interrompues 36~heures.
\end{itemize}

\medskip
\textbf{Analyse} (confiance élevée). Le RTO des expéditions, fixé à
48~heures, est tenu. Les rayonnages surélevés, prévus pour septembre
(Chapitre~\ref{chap:traiter}), n'étaient pas encore posés~: les colis au
sol sont ceux qui ont été endommagés.
\end{boxscenario}

% ============================================================================
\section{Les Trois Plans}
\label{sec:c18-plans}
% ============================================================================

\begin{tableaucours}{L{3.6cm}L{4.2cm}L{4.4cm}}{%
    \tableauHeader{Plan} &
    \tableauHeader{Question} &
    \tableauHeader{Responsable chez LiZbiZ}
}
Plan de gestion de crise & Qui décide, qui communique~? & M.~ETEME (DG) \\
Plan de continuité & Comment continuer l'activité~? & Chaque directeur,
pour son activité \\
Plan de reprise informatique & Comment rétablir les systèmes~? &
M.~KOUNOU \\
\end{tableaucours}
\vspace{-0.8em}
\begin{center}{\small\itshape Tableau~18.3. Trois plans, trois questions}\end{center}

\begin{boxobjet}{R-11 · Inondation de Bonabéri (après MOUNGO)}
\textbf{Scénario~:} inondation de l'entrepôt en saison des pluies.
\textbf{Cotation~:} V passe de 2 à 3 (événement survenu, saisons plus
pluvieuses), I reste à 3. \textbf{Traitement~:} plan de continuité (tenu),
rayonnages surélevés (à accélérer), assurance des marchandises (à
vérifier). \textbf{Propriétaire~:} M.~NGUELE. \textbf{Indicateur~:} part des
colis stockés au sol entre juin et octobre.
\end{boxobjet}

\begin{boxmetier}{Responsable de la continuité}
\textbf{Imprimer le plan.} Le responsable de la continuité garde une copie
papier du plan et des contacts dans chaque site et chez chaque suppléant~:
un plan stocké sur le serveur qui vient de tomber n'existe pas.
\end{boxmetier}

\separateur

\begin{boxresume}
\textbf{Ce qu'il faut retenir de ce chapitre~:}
\begin{itemize}
    \item \textbf{L'ISO~22301} est certifiable~; sa clause~8 organise
          analyse, stratégies, plans et exercices.
    \item \textbf{Une stratégie par ressource}~: personnes, locaux,
          technologies, information, fournisseurs.
    \item \textbf{Le mode dégradé} utilise des ressources indépendantes,
          préparées et entraînées.
    \item \textbf{Le plan} dit qui déclenche, qui fait quoi, avec quoi,
          jusqu'à quand.
    \item \textbf{Crise, continuité, reprise}~: trois plans articulés.
\end{itemize}
\end{boxresume}

\begin{boxinfo}
\textbf{Points de contrôle de l'élève.} Je sais~: situer l'ISO~22301 et sa
clause~8~; choisir une stratégie par ressource~; concevoir un mode
dégradé~; structurer un plan de continuité~; distinguer les trois plans.
\end{boxinfo}

\bigskip

\begin{boxexo}[title={\ding{45}\quad Exercice~18.1. Ressources \niveaubase}]
Rangez dans les cinq types de ressources~: le gardien de Bonabéri~;
l'entrepôt de Bassa~; les carnets de bordereaux~; le second opérateur.
\end{boxexo}

\begin{boxexo}[title={\ding{45}\quad Exercice~18.2. Le RTO tenu \niveaubase}]
Calculez la durée d'interruption des expéditions lors de MOUNGO, du
signalement à la reprise. Le RTO de 48~heures est-il tenu~?
\end{boxexo}

\begin{boxexo}[title={\ding{45}\quad Exercice~18.3. Quel plan~? \niveaubase}]
Quel plan répond à~: « qui appelle les clients prioritaires~? »~;
« comment restaurer la base de facturation~? »~; « faut-il informer la
presse~? »
\end{boxexo}

\begin{boxexo}[title={\ding{45}\quad Exercice~18.4. Indépendance \niveaumoyen}]
Vérifiez que chaque mode dégradé du Tableau~18.2 est indépendant des
ressources qu'il remplace.
\end{boxexo}

\begin{boxexo}[title={\ding{45}\quad Exercice~18.5. Déclenchement \niveaumoyen}]
Rédigez les critères de déclenchement du plan de continuité des
expéditions de Douala.
\end{boxexo}

\begin{boxexo}[title={\ding{45}\quad Exercice~18.6. L'accord de repli \niveaumoyen}]
Quelles clauses l'accord avec l'entrepôt partenaire de Bassa doit-il
contenir~?
\end{boxexo}

\begin{boxexo}[title={\ding{45}\quad Exercice~18.7. Les colis au sol \niveaumoyen}]
Les colis endommagés étaient ceux stockés au sol. Quelle leçon pour le
plan de traitement et pour l'indicateur de R-11~?
\end{boxexo}

\begin{boxexo}[title={\ding{45}\quad Exercice~18.8. Fin de crise \niveaumoyen}]
Proposez les critères de retour à la normale après MOUNGO.
\end{boxexo}

\begin{boxexo}[title={\ding{45}\quad Exercice~18.9. Un plan complet \niveauexpert}]
Rédigez le plan de continuité de la facturation de LiZbiZ, selon la
structure en sept points.
\end{boxexo}

\begin{boxexo}[title={\ding{45}\quad Exercice~18.10. Stratégie bancaire \niveauexpert}]
À partir de l'analyse d'impact du Chapitre~\ref{chap:bia}, proposez à la
BANQUE SANAGA une stratégie pour chacune des cinq ressources.
\end{boxexo}

\bigskip

\begin{boxinfo}
\textbf{Pour aller plus loin.}
\begin{itemize}
    \item ISO~22301~\cite{iso22301_2019}~; ISO~22313~\cite{iso22313_2020}.
\end{itemize}
\end{boxinfo}

\begin{boxinfo}
\textbf{Transition.} Le plan de reprise informatique est la partie la plus
technique. Le Chapitre~\ref{chap:pra} la traite, avec la résilience
opérationnelle.
\end{boxinfo}
